\subsubsection*{\texorpdfstring{Majorana \(\{+,-\}\) proper subgroups}
{Majorana \{+,-\} proper subgroups}}

The Majorana class requires two physical time reversals whose squares have
opposite signs and which anticommute.  In the formulas below, the first tensor
factor acts on the Nambu components and the second acts on a two-component
internal space.  We write \(I=I_2\).  A factor \(\ii\) outside square brackets
multiplies the complete tensor product inside the brackets.  For each case,
\(\Theta_+\) and \(\Theta_-\) are the two transformations that fix the retained
generators and obey
\begin{equation}
    \Theta_+^2=+I,\qquad
    \Theta_-^2=-I,\qquad
    \Theta_+\Theta_-=-\Theta_-\Theta_+.
    \label{eq:proper_majorana_pair}
\end{equation}
The added transformation can be either member of this pair.  Thus
\(\Theta_3^2\) is positive in some cases and negative in others.

For every case, the generator space displayed below is the complete solution
of Eqs.~\eqref{eq:group_only_generators} and
\eqref{eq:extra_TR_restriction} in the stated four-component representation.

\medskip\noindent\textit{Bosons.}\quad
Cases B01, B03, and B09 retain the same three generator directions,
\begin{equation*}
 A_{\rm group}\in\operatorname{span}_{\mathbb R}
 \left\{
 \sigma_x\otimes I,\ 
 (\ii\sigma_y)\otimes I,\ 
 \sigma_z\otimes I
 \right\}.
\end{equation*}
Their matrices and the Majorana pairs are
\begin{align*}
\mathrm{B01}:\quad
 &(J_1,J_2,U_3)=
 \bigl(I\otimes I,\,I\otimes\sigma_x,\,
 I\otimes(\ii\sigma_y)\bigr),\\[-2pt]
 &(\Theta_+,\Theta_-)=
 \bigl((J_2\mc K),(U_3\mc K)\bigr)
 =\bigl((I\otimes\sigma_x)\mc K,
 [I\otimes(\ii\sigma_y)]\mc K\bigr),\\[4pt]
\mathrm{B03}:\quad
 &(J_1,J_2,U_3)=
 \bigl(I\otimes I,\,I\otimes(\ii\sigma_y),\,
 I\otimes\sigma_x\bigr),\\[-2pt]
 &(\Theta_+,\Theta_-)=
 \bigl((U_3\mc K),(J_2\mc K)\bigr)
 =\bigl((I\otimes\sigma_x)\mc K,
 [I\otimes(\ii\sigma_y)]\mc K\bigr),\\[4pt]
\mathrm{B09}:\quad
 &(J_1,J_2,U_3)=
 \bigl(I\otimes(\ii\sigma_y),\,I\otimes I,\,
 I\otimes\sigma_x\bigr),\\[-2pt]
 &(\Theta_+,\Theta_-)=
 \bigl((U_3\mc K),(J_1\mc K)\bigr)
 =\bigl((I\otimes\sigma_x)\mc K,
 [I\otimes(\ii\sigma_y)]\mc K\bigr).
\end{align*}

For B06, the complete retained space is
\begin{equation*}
 A_{\rm group}\in\operatorname{span}_{\mathbb R}
 \left\{
 \ii[\sigma_x\otimes\sigma_z],\ 
 \ii[(\ii\sigma_y)\otimes\sigma_z],\ 
 \sigma_z\otimes I
 \right\},
\end{equation*}
with
\begin{align*}
\mathrm{B06}:\quad
 &(J_1,J_2,U_3)=
 \bigl(I\otimes\sigma_x,\,\sigma_z\otimes\sigma_z,\,
 I\otimes(\ii\sigma_y)\bigr),\\[-2pt]
 &(\Theta_+,\Theta_-)=
 \bigl((J_1\mc K),(U_3\mc K)\bigr)
 =\bigl((I\otimes\sigma_x)\mc K,
 [I\otimes(\ii\sigma_y)]\mc K\bigr).
\end{align*}

Cases B08 and B26 retain
\begin{equation*}
 A_{\rm group}\in\operatorname{span}_{\mathbb R}
 \left\{
 \sigma_x\otimes\sigma_x,\ 
 (\ii\sigma_y)\otimes\sigma_x,\ 
 \sigma_z\otimes I
 \right\}.
\end{equation*}
Their explicit realizations are
\begin{align*}
\mathrm{B08}:\quad
 &(J_1,J_2,U_3)=
 \bigl(I\otimes\sigma_x,\,
 \sigma_z\otimes(\ii\sigma_y),\,I\otimes I\bigr),\\[-2pt]
 &(\Theta_+,\Theta_-)=
 \bigl((U_3\mc K),(\ii J_2\mc K)\bigr)
 =\bigl((I\otimes I)\mc K,\,
 \ii[\sigma_z\otimes(\ii\sigma_y)]\mc K\bigr),\\[4pt]
\mathrm{B26}:\quad
 &(J_1,J_2,U_3)=
 \bigl(\sigma_z\otimes(\ii\sigma_y),\,
 I\otimes\sigma_x,\,I\otimes I\bigr),\\[-2pt]
 &(\Theta_+,\Theta_-)=
 \bigl((U_3\mc K),(\ii J_1\mc K)\bigr)
 =\bigl((I\otimes I)\mc K,\,
 \ii[\sigma_z\otimes(\ii\sigma_y)]\mc K\bigr).
\end{align*}

Cases B14, B18, B20, and B21 share the complete retained space
\begin{equation*}
 A_{\rm group}\in\operatorname{span}_{\mathbb R}
 \left\{
 \ii[\sigma_x\otimes\sigma_x],\ 
 \ii[(\ii\sigma_y)\otimes\sigma_x],\ 
 \sigma_z\otimes I
 \right\}.
\end{equation*}
The transformations are
\begin{align*}
\mathrm{B14}:\quad
 &(J_1,J_2,U_3)=
 \bigl(I\otimes(\ii\sigma_y),\,
 \sigma_z\otimes\sigma_x,\,I\otimes\sigma_z\bigr),\\[-2pt]
 &(\Theta_+,\Theta_-)=
 \bigl((U_3\mc K),(J_1\mc K)\bigr)
 =\bigl((I\otimes\sigma_z)\mc K,\,
 [I\otimes(\ii\sigma_y)]\mc K\bigr),\\[4pt]
\mathrm{B18}:\quad
 &(J_1,J_2,U_3)=
 \bigl(\sigma_z\otimes\sigma_x,\,I\otimes\sigma_z,\,
 I\otimes(\ii\sigma_y)\bigr),\\[-2pt]
 &(\Theta_+,\Theta_-)=
 \bigl((J_2\mc K),(U_3\mc K)\bigr)
 =\bigl((I\otimes\sigma_z)\mc K,\,
 [I\otimes(\ii\sigma_y)]\mc K\bigr),\\[4pt]
\mathrm{B20}:\quad
 &(J_1,J_2,U_3)=
 \bigl(\sigma_z\otimes\sigma_x,\,
 I\otimes(\ii\sigma_y),\,I\otimes\sigma_z\bigr),\\[-2pt]
 &(\Theta_+,\Theta_-)=
 \bigl((U_3\mc K),(J_2\mc K)\bigr)
 =\bigl((I\otimes\sigma_z)\mc K,\,
 [I\otimes(\ii\sigma_y)]\mc K\bigr),\\[4pt]
\mathrm{B21}:\quad
 &(J_1,J_2,U_3)=
 \bigl(\sigma_z\otimes I,\,\sigma_z\otimes\sigma_x,\,
 I\otimes(\ii\sigma_y)\bigr),\\[-2pt]
 &(\Theta_+,\Theta_-)=
 \bigl((\ii J_1\mc K),(U_3\mc K)\bigr)
 =\bigl(\ii(\sigma_z\otimes I)\mc K,\,
 [I\otimes(\ii\sigma_y)]\mc K\bigr).
\end{align*}

Finally, B23 and B29 retain
\begin{equation*}
 A_{\rm group}\in\operatorname{span}_{\mathbb R}
 \left\{
 \ii[\sigma_x\otimes I],\ 
 \ii[(\ii\sigma_y)\otimes I],\ 
 \sigma_z\otimes I
 \right\}.
\end{equation*}
Their matrices and Majorana pairs are
\begin{align*}
\mathrm{B23}:\quad
 &(J_1,J_2,U_3)=
 \bigl(\sigma_z\otimes I,\,
 \sigma_z\otimes(\ii\sigma_y),\,
 \ii[\sigma_z\otimes\sigma_x]\bigr),\\[-2pt]
 &(\Theta_+,\Theta_-)=
 \bigl((U_3\mc K),(\ii J_2\mc K)\bigr)
 =\bigl(\ii(\sigma_z\otimes\sigma_x)\mc K,\,
 \ii[\sigma_z\otimes(\ii\sigma_y)]\mc K\bigr),\\[4pt]
\mathrm{B29}:\quad
 &(J_1,J_2,U_3)=
 \bigl(\sigma_z\otimes(\ii\sigma_y),\,
 \sigma_z\otimes I,\,
 \ii[\sigma_z\otimes\sigma_x]\bigr),\\[-2pt]
 &(\Theta_+,\Theta_-)=
 \bigl((U_3\mc K),(\ii J_1\mc K)\bigr)
 =\bigl(\ii(\sigma_z\otimes\sigma_x)\mc K,\,
 \ii[\sigma_z\otimes(\ii\sigma_y)]\mc K\bigr).
\end{align*}
In each bosonic case the displayed space is a flavor-resolved realization of
the three quadratic directions
formed from a number-conserving term and two independent
particle-number-changing terms.

\medskip\noindent\textit{Fermions.}\quad
For F01 the complete retained space is
\(A_{\rm group}\in\operatorname{span}_{\mathbb R}
\{\sigma_z\otimes I\}\), with
\begin{align*}
\mathrm{F01}:\quad
 &(J_1,J_2,U_3)=
 \bigl(I\otimes I,\,I\otimes\sigma_x,\,
 I\otimes(\ii\sigma_y)\bigr),\\[-2pt]
 &(\Theta_+,\Theta_-)=
 \bigl((J_2\mc K),(U_3\mc K)\bigr)
 =\bigl((I\otimes\sigma_x)\mc K,\,
 [I\otimes(\ii\sigma_y)]\mc K\bigr).
\end{align*}

Cases F02 and F08 retain
\(A_{\rm group}\in\operatorname{span}_{\mathbb R}
\{\ii[I\otimes(\ii\sigma_y)],
\ii[(\ii\sigma_y)\otimes(\ii\sigma_y)]\}\).
Their explicit transformations are
\begin{align*}
\mathrm{F02}:\quad
 &(J_1,J_2,U_3)=
 \bigl(I\otimes\sigma_x,\,I\otimes\sigma_z,\,
 \ii[(\ii\sigma_y)\otimes\sigma_x]\bigr),\\[-2pt]
 &(\Theta_+,\Theta_-)=
 \bigl((J_1\mc K),(U_3\mc K)\bigr)
 =\bigl((I\otimes\sigma_x)\mc K,\,
 \ii[(\ii\sigma_y)\otimes\sigma_x]\mc K\bigr),\\[4pt]
\mathrm{F08}:\quad
 &(J_1,J_2,U_3)=
 \bigl(I\otimes\sigma_x,\,
 (\ii\sigma_y)\otimes\sigma_z,\,I\otimes\sigma_z\bigr),\\[-2pt]
 &(\Theta_+,\Theta_-)=
 \bigl((U_3\mc K),(\ii J_2\mc K)\bigr)
 =\bigl((I\otimes\sigma_z)\mc K,\,
 \ii[(\ii\sigma_y)\otimes\sigma_z]\mc K\bigr).
\end{align*}

Cases F03, F05, F09, F17, F23, and F29 share
\(A_{\rm group}\in\operatorname{span}_{\mathbb R}
\{I\otimes(\ii\sigma_y),
(\ii\sigma_y)\otimes(\ii\sigma_y)\}\).
The six realizations are
\begin{align*}
\mathrm{F03}:\quad
 &(J_1,J_2,U_3)=
 \bigl(I\otimes I,\,I\otimes(\ii\sigma_y),\,
 \ii[(\ii\sigma_y)\otimes I]\bigr),\\[-2pt]
 &(\Theta_+,\Theta_-)=
 \bigl((J_1\mc K),(U_3\mc K)\bigr)
 =\bigl((I\otimes I)\mc K,\,
 \ii[(\ii\sigma_y)\otimes I]\mc K\bigr),\\[4pt]
\mathrm{F05}:\quad
 &(J_1,J_2,U_3)=
 \bigl(I\otimes I,\,
 (\ii\sigma_y)\otimes(\ii\sigma_y),\,
 I\otimes(\ii\sigma_y)\bigr),\\[-2pt]
 &(\Theta_+,\Theta_-)=
 \bigl((\ii J_2\mc K),(U_3\mc K)\bigr)
 =\bigl(\ii[(\ii\sigma_y)\otimes(\ii\sigma_y)]\mc K,\,
 [I\otimes(\ii\sigma_y)]\mc K\bigr),\\[4pt]
\mathrm{F09}:\quad
 &(J_1,J_2,U_3)=
 \bigl(I\otimes(\ii\sigma_y),\,I\otimes I,\,
 \ii[(\ii\sigma_y)\otimes I]\bigr),\\*[-2pt]
 &(\Theta_+,\Theta_-)=
 \bigl((J_2\mc K),(U_3\mc K)\bigr)
 =\bigl((I\otimes I)\mc K,\,
 \ii[(\ii\sigma_y)\otimes I]\mc K\bigr),\\[4pt]
\mathrm{F17}:\quad
 &(J_1,J_2,U_3)=
 \bigl((\ii\sigma_y)\otimes(\ii\sigma_y),\,
 I\otimes I,\,I\otimes(\ii\sigma_y)\bigr),\\[-2pt]
 &(\Theta_+,\Theta_-)=
 \bigl((\ii J_1\mc K),(U_3\mc K)\bigr)
 =\bigl(\ii[(\ii\sigma_y)\otimes(\ii\sigma_y)]\mc K,\,
 [I\otimes(\ii\sigma_y)]\mc K\bigr),\\[4pt]
\mathrm{F23}:\quad
 &(J_1,J_2,U_3)=
 \bigl((\ii\sigma_y)\otimes(\ii\sigma_y),\,
 (\ii\sigma_y)\otimes I,\,I\otimes I\bigr),\\[-2pt]
 &(\Theta_+,\Theta_-)=
 \bigl((U_3\mc K),(\ii J_2\mc K)\bigr)
 =\bigl((I\otimes I)\mc K,\,
 \ii[(\ii\sigma_y)\otimes I]\mc K\bigr),\\[4pt]
\mathrm{F29}:\quad
 &(J_1,J_2,U_3)=
 \bigl((\ii\sigma_y)\otimes I,\,
 (\ii\sigma_y)\otimes(\ii\sigma_y),\,I\otimes I\bigr),\\[-2pt]
 &(\Theta_+,\Theta_-)=
 \bigl((U_3\mc K),(\ii J_1\mc K)\bigr)
 =\bigl((I\otimes I)\mc K,\,
 \ii[(\ii\sigma_y)\otimes I]\mc K\bigr).
\end{align*}

Cases F06 and F18 retain
\(A_{\rm group}\in\operatorname{span}_{\mathbb R}
\{\sigma_z\otimes\sigma_x\}\).
Their transformations are
\begin{align*}
\mathrm{F06}:\quad
 &(J_1,J_2,U_3)=
 \bigl(I\otimes\sigma_x,\,
 (\ii\sigma_y)\otimes(\ii\sigma_y),\,
 \sigma_x\otimes(\ii\sigma_y)\bigr),\\[-2pt]
 &(\Theta_+,\Theta_-)=
 \bigl((J_1\mc K),(U_3\mc K)\bigr)
 =\bigl((I\otimes\sigma_x)\mc K,\,
 [\sigma_x\otimes(\ii\sigma_y)]\mc K\bigr),\\[4pt]
\mathrm{F18}:\quad
 &(J_1,J_2,U_3)=
 \bigl((\ii\sigma_y)\otimes(\ii\sigma_y),\,
 I\otimes\sigma_x,\,
 \sigma_x\otimes(\ii\sigma_y)\bigr),\\[-2pt]
 &(\Theta_+,\Theta_-)=
 \bigl((J_2\mc K),(U_3\mc K)\bigr)
 =\bigl((I\otimes\sigma_x)\mc K,\,
 [\sigma_x\otimes(\ii\sigma_y)]\mc K\bigr).
\end{align*}

Cases F14, F20, F21, F22, and F26 have distinct retained spaces.  For F14,
\(A_{\rm group}\in\operatorname{span}_{\mathbb R}
\{\sigma_z\otimes I,\ \ii[\sigma_z\otimes\sigma_z]\}\), with
\begin{align*}
\mathrm{F14}:\quad
 &(J_1,J_2,U_3)=
 \bigl(I\otimes(\ii\sigma_y),\,
 \sigma_z\otimes\sigma_x,\,I\otimes\sigma_x\bigr),\\[-2pt]
 &(\Theta_+,\Theta_-)=
 \bigl((U_3\mc K),(J_1\mc K)\bigr)
 =\bigl((I\otimes\sigma_x)\mc K,\,
 [I\otimes(\ii\sigma_y)]\mc K\bigr).
\end{align*}
For F20,
\(A_{\rm group}\in\operatorname{span}_{\mathbb R}
\{\ii[\sigma_z\otimes I],\ \sigma_z\otimes\sigma_z\}\), with
\begin{align*}
\mathrm{F20}:\quad
 &(J_1,J_2,U_3)=
 \bigl((\ii\sigma_y)\otimes(\ii\sigma_y),\,
 \sigma_x\otimes(\ii\sigma_y),\,
 \sigma_x\otimes\sigma_x\bigr),\\[-2pt]
 &(\Theta_+,\Theta_-)=
 \bigl((U_3\mc K),(J_2\mc K)\bigr)
 =\bigl((\sigma_x\otimes\sigma_x)\mc K,\,
 [\sigma_x\otimes(\ii\sigma_y)]\mc K\bigr).
\end{align*}
For F21,
\(A_{\rm group}\in\operatorname{span}_{\mathbb R}
\{(\ii\sigma_y)\otimes(\ii\sigma_y)\}\), with
\begin{align*}
\mathrm{F21}:\quad
 &(J_1,J_2,U_3)=
 \bigl((\ii\sigma_y)\otimes(\ii\sigma_y),\,
 \sigma_z\otimes\sigma_x,\,I\otimes(\ii\sigma_y)\bigr),\\[-2pt]
 &(\Theta_+,\Theta_-)=
 \bigl((\ii J_1\mc K),(U_3\mc K)\bigr)
 =\bigl(\ii[(\ii\sigma_y)\otimes(\ii\sigma_y)]\mc K,\,
 [I\otimes(\ii\sigma_y)]\mc K\bigr).
\end{align*}
For F22,
\(A_{\rm group}\in\operatorname{span}_{\mathbb R}
\{I\otimes(\ii\sigma_y),
\ii[\sigma_x\otimes(\ii\sigma_y)]\}\), with
\begin{align*}
\mathrm{F22}:\quad
 &(J_1,J_2,U_3)=
 \bigl((\ii\sigma_y)\otimes(\ii\sigma_y),\,
 \sigma_z\otimes I,\,
 \ii[(\ii\sigma_y)\otimes I]\bigr),\\[-2pt]
 &(\Theta_+,\Theta_-)=
 \bigl((\ii J_2\mc K),(U_3\mc K)\bigr)
 =\bigl(\ii(\sigma_z\otimes I)\mc K,\,
 \ii[(\ii\sigma_y)\otimes I]\mc K\bigr).
\end{align*}
For F26,
\(A_{\rm group}\in\operatorname{span}_{\mathbb R}
\{(\ii\sigma_y)\otimes(\ii\sigma_y),
\ii[\sigma_z\otimes\sigma_x]\}\), with
\begin{align*}
\mathrm{F26}:\quad
 &(J_1,J_2,U_3)=
 \bigl((\ii\sigma_y)\otimes I,\,
 \sigma_x\otimes\sigma_x,\,
 I\otimes(\ii\sigma_y)\bigr),\\[-2pt]
 &(\Theta_+,\Theta_-)=
 \bigl((J_2\mc K),(U_3\mc K)\bigr)
 =\bigl((\sigma_x\otimes\sigma_x)\mc K,\,
 [I\otimes(\ii\sigma_y)]\mc K\bigr).
\end{align*}
The retained fermionic spaces contain one or two independent directions.
Equation~\eqref{eq:proper_majorana_pair} fixes the Gaussian weight on every
resulting proper subgroup.
